\section{Uniform rescaling of the Dirichlet--Neumann operator}
\label{app:rescaling}

This appendix records the exact change of variables used to map legacy data
from an original spatial period \(L_{\mathrm{old}}\) to \(L=2\pi\). Set
\(a=L_{\mathrm{old}}/L>0\), so \(L_{\mathrm{old}}=aL\). We keep the period in
the notation here: $\G^{(L)}(\eta;h)$ denotes the DNO on $\T_L$ at depth $h$,
with the unnormalized-normal convention of \eqref{eq:dno-definition}.

\subsection{Covariance of the elliptic boundary-value problem}

Define the pullback
\begin{equation}
  (\mathcal P_a f)(X)=f(aX),
  \qquad \mathcal P_a:C_{\mathrm{per}}(\T_{aL})
  \longrightarrow C_{\mathrm{per}}(\T_L).
  \label{eq:rescaling-pullback}
\end{equation}
Here \(C_{\mathrm{per}}(\T_\ell)\) denotes the periodic functions on a
one-dimensional torus of period \(\ell\).
Let $\eta$ and $\xi$ be posed on $\T_{aL}$, let the original depth be
$h_{\mathrm{old}}$, and let $\Phi$ solve the corresponding harmonic-extension
problem.  Under
\begin{equation}
  x=aX,\qquad y=aY,\qquad
  \phi(X,Y)=\Phi(aX,aY),
  \label{eq:rescaling-coordinates}
\end{equation}
the free surface and depth become
\begin{equation}
  \eta^\flat(X)=a^{-1}(\mathcal P_a\eta)(X),
  \qquad h^\flat=a^{-1}h_{\mathrm{old}}.
  \label{eq:rescaling-geometry}
\end{equation}
The chain rule gives
\[
  \phi_{XX}+\phi_{YY}
  =a^2(\Phi_{xx}+\Phi_{yy})(aX,aY)=0,
  \qquad
  \phi_Y(X,-h^\flat)=a\Phi_y(aX,-h_{\mathrm{old}})=0,
\]
and the transformed Dirichlet condition is
\[
  \phi(X,\eta^\flat(X))=(\mathcal P_a\xi)(X).
\]
Thus $\phi$ is precisely the harmonic extension associated with
$\G^{(L)}(\eta^\flat;h^\flat)(\mathcal P_a\xi)$.  Moreover,
\begin{equation}
  \partial_X\eta^\flat=\mathcal P_a\eta_x,
  \qquad
  \Phi_x(aX,aY)=a^{-1}\phi_X(X,Y),
  \qquad
  \Phi_y(aX,aY)=a^{-1}\phi_Y(X,Y).
  \label{eq:rescaling-chain-rule}
\end{equation}
Substitution into the unnormalized normal derivative proves
\begin{equation}
  \boxed{
  \mathcal P_a\!\left[
    \G^{(aL)}(\eta;h_{\mathrm{old}})\xi
  \right]
  =\frac1a\,
  \G^{(L)}\!\left(
    \frac1a\mathcal P_a\eta;\frac{h_{\mathrm{old}}}{a}
  \right)(\mathcal P_a\xi).}
  \label{eq:dno-rescaling}
\end{equation}
The factor $a^{-1}$ comes from the derivative, not from a normalization of the
normal vector. As a consistency check, when $\eta=0$, every periodic function
\(f\) satisfies
\[
  \begin{aligned}
  \mathcal P_a\!\left[
    |D_x|\tanh(h_{\mathrm{old}}|D_x|)f
  \right]
  &=a^{-1}|D_X|\tanh((h_{\mathrm{old}}/a)|D_X|)
    (\mathcal P_a f),\\
  D_x&=-i\partial_x=a^{-1}D_X,
  \qquad D_X=-i\partial_X.
  \end{aligned}
\]

\subsection{Fixed-gravity Zakharov scaling}

Equation \eqref{eq:dno-rescaling} leaves the pulled-back Dirichlet datum
$\mathcal P_a\xi$ unscaled.  To obtain a change of variables for the full
fixed-$g$ evolution, the amplitude of $\xi$ and the time coordinate must also
be rescaled.  Let
$q=\G^{(aL)}(\eta;h_{\mathrm{old}})\xi$, set $t=\sqrt a\,\tau$, and define
\begin{equation}
  \begin{aligned}
    \eta^\sharp(X,\tau)
      &=a^{-1}\eta(aX,\sqrt a\,\tau),
    &\xi^\sharp(X,\tau)
      &=a^{-3/2}\xi(aX,\sqrt a\,\tau),\\
    q^\sharp(X,\tau)
      &=a^{-1/2}q(aX,\sqrt a\,\tau),
    &h^\sharp&=h_{\mathrm{old}}/a.
  \end{aligned}
  \label{eq:zakharov-rescaling}
\end{equation}
Linearity in $\xi$ and \eqref{eq:dno-rescaling} imply
\begin{equation}
  q^\sharp=\G^{(L)}(\eta^\sharp;h^\sharp)\xi^\sharp.
  \label{eq:rescaled-q-identity}
\end{equation}
The first Zakharov equation follows immediately:
\[
  \partial_\tau\eta^\sharp
  =a^{-1/2}\mathcal P_a\eta_t=q^\sharp.
\]
For the second equation, the required factors are
\[
  \partial_\tau\xi^\sharp=a^{-1}\mathcal P_a\xi_t,
  \qquad
  \partial_X\eta^\sharp=\mathcal P_a\eta_x,
  \qquad
  \partial_X\xi^\sharp=a^{-1/2}\mathcal P_a\xi_x.
\]
Consequently every term on the right-hand side of
\eqref{eq:zakharov-xi} acquires the common factor $a^{-1}$:
\begin{align*}
  -g\eta^\sharp&=a^{-1}\mathcal P_a(-g\eta),\\
  -\tfrac12(\xi_X^\sharp)^2
    &=a^{-1}\mathcal P_a(-\tfrac12\xi_x^2),\\
  \frac{(q^\sharp+\eta_X^\sharp\xi_X^\sharp)^2}
       {2(1+(\eta_X^\sharp)^2)}
    &=a^{-1}\mathcal P_a
      \left[\frac{(q+\eta_x\xi_x)^2}{2(1+\eta_x^2)}\right].
\end{align*}
Hence $(\eta^\sharp,\xi^\sharp)$ solves the same fixed-gravity Zakharov
system on $\T_L$ at depth $h^\sharp$.

\subsection{Scope of the identity}

The covariance requires a common dilation of the horizontal and vertical
coordinates; it does not justify changing period and depth independently.
It is also an identity of the continuum DNO and the reference dynamics, not a
claim that a learned residual is exactly equivariant. The analytic
$\G_0+\G_1$ backbone inherits the covariance, whereas the learned residual
used in the reported model also depends on absolute wavenumber, depth, and
dataset-level feature normalization. The rejected residual-only
reparameterization discussed in \cref{sec:rejected-remedies} therefore tests
a model-design choice and does not test \eqref{eq:dno-rescaling} itself.
